\songtitle{Mis brazos, no a mí}

\tag*{2026}\tag{metal-ballad}\tag{spanish-lyrics}\tag{original-source}\tag{thomas-entourage}\tag{thomas-garcia}\tag{sylvia-lai}\label{sec:mis-brazos-no-a-mi}
\begin{notes}
Lyrics by Claude following specific instructions, based on \emph{Shoulders} by Carlos Thompson, from the \emph{Thomas' Entourage} narrative universe, in the voice of Thomas -- Barcelona, November 2026: Sylvia, grown up and finding her place, still runs to his arms, and he wonders what the embrace means. From the same story as \emph{When It Falls Apart} (page \pageref{sec:when-it-falls-apart}), \emph{Her Own Place} (page \pageref{sec:her-own-place}) and \emph{Not a Messenger} (page \pageref{sec:not-a-messenger}). Retrieved from Wyomee Dank's Suno page (V6-MINI, 3 November 2026).
\end{notes}
\begin{style}
Soft metal ballad, unhurried brushed sway; brushed percussion, finger-picked electric guitars and sustained distorted bass; Colombian baritone singer with clear pronunciation, hushed ironic, amused delivery and doubled hooks; close-mic intimacy, a restrained final-chorus lift.
\end{style}
\begin{lyrics}
\songpart{Verse 1}
Barcelona en noviembre, no esperaba aquel abrazo,\\
creí que al hacerse adulta marcaría más el espacio.\\
Pero llegó sin anunciarse, como en otros tiempos,\\
buscando entre mis brazos refugio contra el viento.

Dieciocho horas de excusas para no dejarme ir,\\
y yo fingiendo resistencia, sin querer resistir.\\
Después, en aquella mesa, la volví a encontrar,\\
con su cabello color teal y un mundo por estrenar.

\songpart{Pre-Chorus}
Xumi me habló de ella tomando un café,\\
de algunos tropiezos, de aprender a estar de pie.\\
Que estaba haciendo las cosas, que sabía responder,\\
que nunca usó su apellido para abrirse el camino.

\songpart{Chorus}
No me busca a mí, busca mis brazos,\\
el lugar donde se borran por un rato los fracasos.\\
No es mi nombre lo que viene a reclamar,\\
es este abrazo que no le voy a negar.

No me busca a mí, busca mis brazos,\\
y me pregunto si algún día entenderé\\
que puede hacerse adulta, encontrar su propio espacio,\\
y aun así volver a mí para descansar.

\songpart{Verse 2}
La vi segura y diligente, ocupando su lugar,\\
ya no era aquella niña que tenía que cuidar.\\
No puedo hablar de orgullo como un padre lo haría,\\
todavía tengo tiempo para saber qué sentiría.

Xumi dijo que aprendía, que todo iba bastante bien,\\
que los días complicados no la hicieron retroceder.\\
Y yo pensé en la distancia que creí que iba a elegir,\\
pero hay cosas que no cambian solo porque hay que crecer.

\songpart{Pre-Chorus}
Quizás no era independencia alejarse de los demás,\\
ni madurar significa no mirar nunca hacia atrás.\\
Quizás hay una manera de encontrar dónde encajar\\
sin tener que renunciar a quien te puede abrazar.

\songpart{Chorus}
No me busca a mí, busca mis brazos,\\
el lugar donde se borran por un rato los fracasos.\\
No es mi nombre lo que viene a reclamar,\\
es este abrazo que no le voy a negar.

No me busca a mí, busca mis brazos,\\
y me pregunto si algún día entenderé\\
que puede hacerse adulta, encontrar su propio espacio,\\
y aun así volver a mí para descansar.

\songpart{Bridge}
No soy su padre, no soy su destino,\\
no tengo que mostrarle cuál camino seguir.\\
Ella tiene sus batallas, sus razones, su camino,\\
y yo tengo estos brazos cuando quiera venir.

Tal vez eso sea todo, tal vez sea suficiente,\\
no pedir que me recuerde como fui,\\
ni confundir el refugio con ser imprescindible,\\
ni creer que busca en ellos algo más de lo que vi.

\songpart{Final Chorus}
No me busca a mí, busca mis brazos,\\
y no voy a preguntarle qué hay detrás del abrazo.\\
Que siga construyendo su lugar,\\
que tenga un mundo propio al que regresar.

No me busca a mí, busca mis brazos,\\
y aunque no sepa qué significa para mí,\\
si todavía encuentra un sitio entre mis brazos,\\
no tengo prisa por dejarla partir.
\end{lyrics}

